% =============================================================================
\section{Conclusões}
\label{sec:conclusoes}
\label{sec:conclusao}

Este trabalho documentou a concepção, o desenvolvimento a partir do zero e a
análise empírica rigorosa de um arcabouço completo de Programação Genética para
Regressão Simbólica. A investigação foi orientada pelas diretrizes da
disciplina e articulada em duas fases complementares.

Na \textbf{Fase~1}, a exploração paramétrica canônica sobre três bases
(\code{synth1}, \code{synth2} e \code{concrete1}) demonstrou que o método de
seleção é o fator de maior impacto estrutural: a $\eps$-lexicase superou
expressivamente o torneio e a roleta ($p = 4.12\times10^{-6}$). Em
contrapartida, a replicação expôs três patologias severas da abordagem
clássica: explosões assintóticas de divisão fora da amostra no \code{synth2}
($415.3$ de erro máximo), subajuste por barreira de escala no \code{concrete1}
($R^2 \approx 0.44$), altas taxas de operadores destrutivos ($>88\%$ dos
cruzamentos degradando o melhor progenitor) e inchaço de árvores com
acúmulo de código neutro.

Na \textbf{Fase~2}, a interpretação desses sintomas sob a ótica da
\emph{fitness landscape} orientou a formulação de seis hipóteses mecanicistas
(H1--H6), validadas por ablação isolada e por empilhamento cumulativo em 1.260
execuções adicionais. As conclusões consolidadas incluem:
\begin{enumerate}
  \item A validação do \textbf{Analytic Quotient (AQ)} como substituto
        definitivo da divisão protegida clássica, eliminando 100\% das
        singularidades fora da amostra no \code{synth2} (queda de erro de
        $93.9\%$);
  \item A comprovação de que o \textbf{Linear Scaling} por mínimos quadrados
        desacopla a busca funcional da barreira de amplitude numérica,
        elevando o $R^2$ do \code{concrete1} para $0.78$ com $30/30$ vitórias
        sobre o estado inicial;
  \item A demonstração de que a combinação de \textbf{rejeição de duplicatas},
        \textbf{parcimônia lexicográfica} e \textbf{mutação mista} restringe
        o \emph{bloat} em até $42\%$ e destrava a busca em problemas
        polinomiais, recuperando com exatidão matemática a equação alvo no
        \code{synth1} com RMSE de $0.0006$;
  \item A caracterização do modelo \textbf{\cfg{Stack~B3}} como a recomendação
        ótima global de compromisso, apresentando liderança em posto médio
        geral e ganhos substanciais de generalização em todos os cenários.
\end{enumerate}

Como trabalhos futuros, sugere-se a incorporação de amostragem adaptativa
de casos de treino (\emph{Down-sampled Lexicase}), a otimização contínua de
constantes por algoritmos baseados em gradiente (ex.: BFGS ou Levenberg-Marquardt)
e a integração de restrições de dimensionalidade e análise física via Gramáticas
Evolutivas (Grammatical Evolution).
